\section{Redução de Expoentes (Teorema de Euler Estendido)}

\textbf{Problema base.} Calcular $a^e \bmod m$ quando o expoente $e$ é enorme (string com $10^5$ dígitos, torre $a^{b^c}$).

\textbf{Ideia.} As potências de $a$ módulo $m$ entram num ciclo de tamanho que divide $\varphi(m)$, e o ciclo começa
até o expoente $\varphi(m)$. Vale para \textbf{quaisquer} $a$ e $m$, mesmo com $\gcd(a, m) > 1$:
\[
\boxed{a^e \equiv \begin{cases}
  a^e, & \text{se } e < \varphi(m) \\[2pt]
  a^{(e \bmod \varphi(m)) + \varphi(m)}, & \text{se } e \ge \varphi(m)
\end{cases} \pmod m}
\]

\begin{itemize}
  \item \textbf{Caso $e < \varphi(m)$:} não reduza o expoente; reduzir antes do ciclo muda o resultado quando $\gcd(a, m) > 1$.
  \item \textbf{Caso $e \ge \varphi(m)$:} troque $e$ por $(e \bmod \varphi(m)) + \varphi(m)$; o $+\varphi(m)$ mantém o expoente dentro do ciclo.
  \item \textbf{Se $\gcd(a, m) = 1$ (Euler):} o ciclo começa em $0$, então basta $a^{e \bmod \varphi(m)}$.
\end{itemize}
